\section{Related Work}\label{sec:related}

\emph{Labels.} Label and flow-construction errors in NIDS benchmarks are well documented for CIC-IDS2017 and CSE-CIC-IDS2018~\cite{engelen2021troubleshooting,liu2022error}, and Flood et al.\ catalogue design smells of benchmark datasets~\cite{flood2024smells}. Garcia and Valeros report that most of the datasets they analysed assign labels ``in files or directory names'' (abstract) and that published datasets ``sometimes contain errors that require significant effort to correct'' (Sec.~2.2)~\cite{garcia2023labeling}. For CIC-IoT2023~\cite{neto2023ciciot2023}, Ramkumar et al.\ observed that the captures mix benign and attack traffic and re-filtered them per device by MAC address, but did not quantify the mixing~\cite{ramkumar2024diagnosing}. Vitorino et al.\ re-extract CIC-IoT2023 flows and assume that each capture contains only its named attack~\cite{vitorino2025revisiting}. IoTGeM reports benign packets inside the attack captures of another IoT dataset~\cite{kostas2024iotgem}. Background traffic in attack captures is therefore not new; we add the background share per capture and the measured cost of file-name labels on two datasets.

\emph{Splits.} Random splits of time-ordered data are a known pitfall in security machine learning~\cite{arp2022dos,pendlebury2019tesseract,apruzzese2023sok}, and tree classifiers degrade over time on MAWI flows~\cite{schraven2026mawiflow}. Temporal evaluations involving CIC-IoT2023 exist: Traffic-MoE uses a per-class chronological split on packet bytes~\cite{zhou2026trafficmoe}, BRIDGE an early/late split in a pooled benchmark~\cite{bridge2026tchnet}, and Rasheed and Alobaedy find that drift causes no consistent decrease in performance but ``specific behaviour changes, such as total blindness, sensitivity inflation and relatively stable detection''~\cite{rasheed2026drift}. IoTGeM trains on telnet and tests on RTSP brute force, by design, and finds that the variants do not generalise~\cite{kostas2024iotgem}. Our split is within each capture and traces the gap to its cause, to boundary-spanning connections and, on IoT-23, to missing classes.

\emph{Live detection and exporters.} Learned detectors on IoT gateways are common~\cite{realguard2022}, as are Mininet/Open vSwitch~\cite{lantz2010mininet,pfaff2015ovs} prototypes that block flagged hosts and combine NFStream~\cite{aouini2022nfstream} with a classifier. Flow export and its timeouts are surveyed by Hofstede et al.~\cite{hofstede2014flow}; Pinto et al.\ show that the choice of exporter changes detection results~\cite{pinto2025flowexporter,pinto2025hera}, and Vitorino et al.\ that it changes what a model sees~\cite{vitorino2025revisiting}. Our gateway is not new in its parts; we contribute the fidelity check and what it exposed.

\emph{Response.} Blocking via OpenFlow/OVS drop rules and gating a block on several verdicts appear in SDN mitigation work~\cite{bawany2017ddos,yan2016sdn}; our ARP rule is the classic arpwatch idea~\cite{arpwatch}. A recent SoK reviews how machine-learning SDN defences are evaluated~\cite{orsdemir2026sok}; few works we found report false-block counts on replayed traffic.
